% SGA 6, Exposición VI: fragmento español alineado con la fuente.
% Alcance: sección 2, números 2.3.3--2.6, líneas 9014--9088 inclusive.
% Fuente francesa: sga6_fr_r1/src/fr.tex.
% SHA-256 de la fuente francesa completa:
% 0D2158B9FE3B466309A92165727A0B191422E05FBAA0DACD5806E8322C54BA1A.
\providecommand{\sisi}[2]{#2}

\textbf{2.3.3.} De 2.3.1 y 2.3.2 resulta que el funtor \(\Proj_0\) es exacto. Volviendo a las notaciones de 2.2, se deduce, pues, pasando a las categorías derivadas, un funtor denotado también por \(\Proj_0\):
\[
  \Proj_0:D^+(\mathcal G)\longrightarrow D^+(\mathcal A).
\]
Se ve igualmente en las construcciones anteriores que el funtor \(\Proj_0\) sobre la categoría de los \(B\)-módulos graduados es compatible con el producto tensorial; lo mismo sigue siendo cierto para el funtor obtenido sobre las categorías derivadas.


\textbf{2.4.} Para todo entero $k$, y todo $\mathcal O_X$-módulo $L$, pongamos:
\begin{equation}\tag{2.4.1}
  F_k(L)=\bigoplus_{n\geq k} f_*\bigl(L(n)\bigr).
\end{equation}
El funtor $F_k:\mathcal A\to\mathcal G$ así definido es exacto por la izquierda. Como la categoría de los $\mathcal O_X$-módulos tiene suficientes inyectivos, se puede definir un funtor derivado:
\begin{equation}\tag{2.4.2}
  \mathbf R F_k:D^+(\mathcal A)\longrightarrow D^+(\mathcal G).
\end{equation}
Si $L$ es un $\mathcal O_X$-módulo localmente libre de tipo finito, existe un entero $k_0$ tal que, para $k\geq k_0$, $L$ es $F_k$-acíclico. Para verlo, basta mostrar que, para $n$ suficientemente grande, $L$ es acíclico para los funtores $L\leadsto f_*\bigl(L(n)\bigr)$. Como la tensorización por $\mathcal O_X(n)$ es un funtor exacto, esto resulta inmediatamente de la sucesión espectral de los funtores compuestos y de 1.3.

\textbf{2.5.} Nos proponemos construir, para todo $\mathcal O_X$-módulo $F$, un homomorfismo funtorial en $F$:
\begin{equation}\tag{2.5.1}
  \beta:\operatorname{Proj}_0\bigl(F_k(F)\bigr)\longrightarrow F.
\end{equation}
Para ello, se define la restricción de $\beta$ a los $D_+(U,s)$ (2.3.2). Ahora bien, por construcción se tiene, poniendo $M=F_k(F)$:
\[
  \operatorname{Proj}_0(M)|D_+(U,s)
  =\bar f^{*}_{U,s}\bigl(M_{(s),U}\bigr).
\]
Para definir un homomorfismo
\[
  \bar f^{*}_{U,s}\bigl(M_{(s),U}\bigr)\longrightarrow F|_{D_+(U,s)},
\]
basta, según la definición de $\bar f^{*}_{U,s}$ como adjunto por la izquierda al funtor de imagen directa, definir un homomorfismo $f_{U,s*}(\mathcal O_X)$-lineal
\[
  M_{(s),U}\longrightarrow f_{U,s*}\bigl(F|_{D_+(U,s)}\bigr).
\]
Para hacerlo, hay que darse, para todo abierto $V$ de $U$, un homomorfismo:
\begin{equation}\tag{2.5.2}
  \Gamma\bigl(V,M\bigr)_{(s)}\longrightarrow \Gamma\bigl(D_+(V,s),F\bigr)
\end{equation}
compatible con las restricciones a los abiertos de $U$, y tomar el homomorfismo correspondiente sobre los haces asociados. El homomorfismo (2.5.2) se define de la manera siguiente: sea $z/s^n\in\Gamma(V,M)_{(s)}$, donde, siendo $s$ de grado $d$, $z\in\Gamma(V,M_{nd})=\Gamma\bigl(D_+(V,s),F(nd)\bigr)$; se asocia entonces a $z/s^n$ la sección $z.\sisi{\alpha_d}{\alpha_{nd}}(s^n)^{-1}$ de $F$ sobre $D_+(V,s)$, denotando por $\sisi{\alpha_d}{\alpha_{nd}}$ el homomorfismo canónico
\[
  \Gamma(V,\sisi{B_d}{B_{nd}})\longrightarrow \Gamma\bigl(f^{-1}(V),\mathcal O_X(\sisi{d}{nd})\bigr)
\]
(cf. EGA II 2.6.3 y 2.6.4). Se verifican sin dificultad las compatibilidades necesarias, de donde se obtiene, por pegado, el homomorfismo (2.5.1).

Está claro que, cuando $F$ es cuasicoherente, el homomorfismo así construido es el de EGA II (3.3.4.1).

El homomorfismo de funtor $\operatorname{Proj}_0\circ F_k\to \operatorname{Id}$ sobre la categoría de los $\mathcal O_X$-módulos da, pasando a las categorías derivadas, un homomorfismo de funtor sobre $D^+(\mathcal A)$:
\begin{equation}\tag{2.5.3}
  \operatorname{Proj}_0\circ\mathbf R F_k\longrightarrow \operatorname{Id}_{D^+(\mathcal A)}.
\end{equation}

\textbf{Proposición 2.6.} Sean $S$ un esquema cuasicompacto, $E$ un $\mathcal O_S$-módulo localmente libre de tipo finito, $X=\mathbf P(E)$, y $L^{\bullet}$ un complejo perfecto sobre $X$. Entonces existe un entero $k_0$ tal que, para $k\geq k_0$, se verifican las propiedades siguientes:
\[
\begin{array}{ll}
\mathrm{i)}& \mathbf R F_k(L^{\bullet})\text{ es }\mathcal G_0\text{-perfecto};\\[2mm]
\mathrm{ii)}& \text{El morfismo canónico }(2.5.3):
\operatorname{Proj}_0\circ\mathbf R F_k(L^{\bullet})\longrightarrow L^{\bullet}
\text{ es un isomorfismo de }D^+(\mathcal A).
\end{array}
\]

La noción de complejo $\mathcal G_0$-perfecto considerada aquí es local, es decir, se trata de perfección respecto de la categoría fibrada sobre $S$ definida por $\mathcal G_0$. Esta noción de perfección tiene sentido, pues $\mathcal G_0$ verifica los axiomas de (I 4.0): solo la estabilidad por núcleo de epimorfismo no es inmediata; resulta de 1.5.1.

Como $S$ es cuasicompacto, basta verificar localmente sobre $S$ las dos afirmaciones de 2.6, de modo que se puede suponer $S$ afín. El haz invertible $\mathcal O_X(1)$ es entonces amplio, y se dispone de resoluciones globales sobre $X$ (II 2.2.9). Existe, pues, un complejo $L'{}^{\bullet}$ de $D_{\mathrm{parf}}(\mathcal A)$, estrictamente perfecto, y un isomorfismo en $D_{\mathrm{parf}}(\mathcal A)$, $L^{\bullet}\xleftarrow{\sim}L'{}^{\bullet}$, lo que permite suponer $L^{\bullet}$ estrictamente perfecto. Ahora bien, se vio en 2.4 que, para $k$ suficientemente grande, los $L^i$, todos nulos salvo un número finito, son $F_k$-acíclicos; se tiene por consiguiente:
\[
  \mathbf R F_k(L^{\bullet})=F_k(L^{\bullet})
\]
si $k$ es suficientemente grande. Ahora bien, según 1.3 y 1.6, se puede encontrar $k_0$ tal que, para $k\geq k_0$, se cumpla la condición anterior y los $F_k(L^i)$ sean $B$-módulos $\mathcal G_0$-perfectos para todo $i$. El complejo $F_k(L^{\bullet})$ es entonces $\mathcal G_0$-perfecto, de donde i).

Si $k$ se elige como anteriormente, el homomorfismo (2.5.3) es entonces el homomorfismo
\[
  \operatorname{Proj}_0\bigl(F_k(L^{\bullet})\bigr)\longrightarrow L^{\bullet},
\]
y es un isomorfismo sobre cada componente, puesto que se ha supuesto que $L^{\bullet}$ es estrictamente perfecto, de donde ii).

